\section{Reproduction}
\label{app:repro}

\paragraph{The measured run.} All results in this report come from the directory
\code{results-0294f37}, produced by a clean Release build of commit \code{0294f37} on
the branch \code{io-variants}. The repository contains, for every cell and repetition,
the metadata (\code{rep*.meta.json}), the generator's result (\code{rep*.stress.json})
and the node's metrics (\code{rep*.metrics.jsonl}), together with \code{summary.csv}
and \code{report.html}. The node and generator logs are not committed.

\paragraph{Build and run.}
\begin{verbatim}
cmake -B build -S . -DCMAKE_BUILD_TYPE=Release && cmake --build build -j
scripts/run_matrix.py --app build/app --stressgen build/stressgen \
    --out results-<commit> --reps 3
scripts/analyze.py results-<commit>      # summary.csv, report.html
\end{verbatim}

\paragraph{Notes.}
\begin{itemize}
  \item The working tree must be clean, so every cell's build stamp names an exact
        commit.
  \item The node runs on CPUs 0--1 and the generator on CPUs 4--7.
  \item An interrupted run resumes when the same command is run again: cells that
        completed are kept and failed ones are retried. It refuses to mix builds.
\end{itemize}

\paragraph{Figures and tables.} Every figure and table in Section~\ref{sec:results}
is generated from \code{results-0294f37} by the script \code{make\_figures.py} in
the directory \code{report/visualisation-scripts}. The handshake diagrams
(Figure~\ref{fig:handshake}) are PlantUML sources in the same directory, rendered by
\code{make\_handshake.sh}.

\section{Test suite}

\begin{itemize}
  \item One unit-test executable per component: queue, wire format, sockets and pools,
        \code{TxPoller}, store, sender, workers, rendering, metrics, \code{stressgen},
        \code{peer-threads} and others.
  \item An end-to-end script that runs two headless nodes in every mode.
  \item Tests for the analysis script.
  \item A normal \code{ctest -j8} passes 26 of 26 tests.
\end{itemize}

\paragraph{Parallel load.} Running six to eight copies of the network-heavy tests at
once still fails about 2\% of runs, because fixed deadlines in the two-node checks
expire under that artificial load (see also Section~\ref{sec:lessons}).

\section{All results}
\label{app:full}

Tables~\ref{tab:full-S1}--\ref{tab:full-S7} list every swept value and every mode of
every scenario. Each value is the median over the valid repetitions; ``valid'' gives
how many of the three repetitions passed the validity checks, and a cell with no
valid repetition shows dashes. Columns:
\begin{itemize}
  \item \textbf{msg/s}: echoed messages per second; \textbf{lost}: messages sent in
        the window and never echoed;
  \item \textbf{p50}, \textbf{p99}, \textbf{p99.9}: round-trip time in ms;
  \item \textbf{CPU \%}: node CPU usage (Equation~\ref{eq:cpu}); \textbf{vol./s} and
        \textbf{invol./s}: voluntary and involuntary context switches per second;
  \item \textbf{thr.} and \textbf{fds}: the node's peak thread and descriptor counts;
  \item \textbf{ready}: peers that completed the handshake; \textbf{hs p50}: median
        handshake time in ms;
  \item \textbf{queued}: sends that had to be queued because the socket's transmit
        buffer was full; \textbf{evict.}: peers evicted by the cap.
\end{itemize}

\begin{landscape}
\scriptsize
\setlength{\tabcolsep}{4pt}
\setlength{\LTcapwidth}{\linewidth}
\begin{longtable}{@{}llrrrrrrrrrrrrrrr@{}}
\caption{S1 (fan-in): all results.}\label{tab:full-S1}\\
\toprule
peers & Variant & valid & msg/s & lost & p50 & p99 & p99.9 & CPU \% & vol./s & invol./s & thr. & fds & ready & hs p50 & queued & evict. \\
\midrule
\endfirsthead
\toprule
peers & Variant & valid & msg/s & lost & p50 & p99 & p99.9 & CPU \% & vol./s & invol./s & thr. & fds & ready & hs p50 & queued & evict. \\
\midrule
\endhead
\bottomrule
\endfoot
1 & \texttt{threads} & 3/3 & 100 & 0 & 0.7 & 1.8 & 2.5 & 2 & 331 & 0 & 9 & 8 & 1 & 1.9 & 0 & 0 \\
 & \texttt{pool-sender} & 3/3 & 100 & 0 & 0.7 & 1.8 & 2.6 & 2 & 335 & 0 & 10 & 9 & 1 & 2.5 & 0 & 0 \\
 & \texttt{peer-threads} & 3/3 & 100 & 0 & 0.7 & 2.1 & 3.1 & 2 & 335 & 0 & 9 & 8 & 1 & 1.4 & 0 & 0 \\
 & \texttt{poll} & 3/3 & 100 & 0 & 0.8 & 1.6 & 2.2 & 2 & 324 & 0 & 8 & 9 & 1 & 2.2 & 0 & 0 \\
 & \texttt{poll-duplex} & 3/3 & 100 & 0 & 0.7 & 1.8 & 2.7 & 2 & 325 & 0 & 8 & 8 & 1 & 1.5 & 0 & 0 \\
 & \texttt{poll-split} & 3/3 & 100 & 0 & 0.7 & 2.2 & 3.2 & 2 & 325 & 0 & 9 & 9 & 1 & 2.2 & 0 & 0 \\
 & \texttt{epoll} & 3/3 & 100 & 0 & 0.7 & 1.9 & 2.9 & 2 & 325 & 0 & 8 & 10 & 1 & 2.0 & 0 & 0 \\
 & \texttt{epoll-duplex} & 3/3 & 100 & 0 & 0.6 & 1.9 & 2.9 & 2 & 325 & 0 & 8 & 9 & 1 & 2.1 & 0 & 0 \\
 & \texttt{epoll-split} & 3/3 & 100 & 0 & 0.7 & 1.9 & 2.8 & 2 & 325 & 0 & 9 & 11 & 1 & 1.8 & 0 & 0 \\
\midrule
8 & \texttt{threads} & 3/3 & 800 & 0 & 0.9 & 1.7 & 2.6 & 4 & 1{,}382 & 55 & 16 & 22 & 8 & 5.1 & 0 & 0 \\
 & \texttt{pool-sender} & 3/3 & 800 & 0 & 0.9 & 1.6 & 2.4 & 4 & 1{,}358 & 60 & 17 & 23 & 8 & 4.6 & 0 & 0 \\
 & \texttt{peer-threads} & 3/3 & 800 & 0 & 1.0 & 2.0 & 2.9 & 3 & 1{,}383 & 53 & 16 & 22 & 8 & 2.9 & 0 & 0 \\
 & \texttt{poll} & 3/3 & 800 & 0 & 1.0 & 1.7 & 2.4 & 5 & 817 & 2 & 8 & 23 & 8 & 4.1 & 0 & 0 \\
 & \texttt{poll-duplex} & 3/3 & 800 & 0 & 0.9 & 1.4 & 2.8 & 4 & 665 & 1 & 8 & 15 & 8 & 2.2 & 0 & 0 \\
 & \texttt{poll-split} & 3/3 & 800 & 0 & 1.0 & 2.2 & 3.4 & 3 & 493 & 3 & 9 & 16 & 8 & 2.2 & 0 & 0 \\
 & \texttt{epoll} & 3/3 & 800 & 0 & 1.0 & 1.7 & 2.4 & 4 & 757 & 1 & 8 & 24 & 8 & 5.0 & 0 & 0 \\
 & \texttt{epoll-duplex} & 3/3 & 800 & 0 & 1.0 & 2.2 & 3.0 & 3 & 542 & 0 & 8 & 16 & 8 & 2.9 & 0 & 0 \\
 & \texttt{epoll-split} & 3/3 & 800 & 0 & 1.0 & 2.2 & 3.3 & 3 & 507 & 3 & 9 & 18 & 8 & 1.8 & 0 & 0 \\
\midrule
64 & \texttt{threads} & 3/3 & 6{,}400 & 0 & 1.8 & 3.8 & 5.5 & 18 & 8{,}409 & 690 & 72 & 134 & 64 & 51.4 & 0 & 0 \\
 & \texttt{pool-sender} & 3/3 & 6{,}400 & 0 & 1.9 & 3.9 & 5.4 & 19 & 8{,}247 & 639 & 73 & 135 & 64 & 46.1 & 0 & 0 \\
 & \texttt{peer-threads} & 3/3 & 6{,}400 & 0 & 1.7 & 3.2 & 3.9 & 17 & 9{,}036 & 786 & 72 & 134 & 64 & 24.6 & 0 & 0 \\
 & \texttt{poll} & 3/3 & 6{,}400 & 0 & 2.4 & 4.5 & 5.2 & 35 & 6{,}672 & 150 & 8 & 135 & 64 & 48.2 & 0 & 0 \\
 & \texttt{poll-duplex} & 3/3 & 6{,}400 & 0 & 2.1 & 3.7 & 4.5 & 24 & 5{,}206 & 146 & 8 & 71 & 64 & 6.7 & 0 & 0 \\
 & \texttt{poll-split} & 3/3 & 6{,}400 & 0 & 1.7 & 3.2 & 4.7 & 20 & 3{,}926 & 1{,}134 & 9 & 72 & 64 & 5.9 & 0 & 0 \\
 & \texttt{epoll} & 3/3 & 6{,}400 & 0 & 2.2 & 3.9 & 4.8 & 26 & 4{,}774 & 102 & 8 & 136 & 64 & 50.3 & 0 & 0 \\
 & \texttt{epoll-duplex} & 3/3 & 6{,}400 & 0 & 2.0 & 3.6 & 4.3 & 19 & 4{,}381 & 107 & 8 & 72 & 64 & 5.6 & 0 & 0 \\
 & \texttt{epoll-split} & 3/3 & 6{,}400 & 0 & 1.7 & 3.0 & 3.7 & 14 & 2{,}472 & 468 & 9 & 74 & 64 & 4.7 & 0 & 0 \\
\midrule
256 & \texttt{threads} & 3/3 & 25{,}600 & 0 & 4.1 & 7.2 & 8.7 & 49 & 28{,}529 & 1{,}308 & 264 & 518 & 256 & 151.0 & 0 & 0 \\
 & \texttt{pool-sender} & 3/3 & 25{,}600 & 0 & 4.0 & 7.1 & 8.9 & 49 & 28{,}615 & 1{,}315 & 265 & 519 & 256 & 142.6 & 0 & 0 \\
 & \texttt{peer-threads} & 3/3 & 25{,}600 & 0 & 3.9 & 6.8 & 9.7 & 54 & 29{,}814 & 1{,}592 & 264 & 518 & 256 & 96.5 & 0 & 0 \\
 & \texttt{poll} & 3/3 & 25{,}600 & 0 & 11.5 & 22.5 & 28.3 & 132 & 25{,}314 & 221 & 8 & 519 & 256 & 117.4 & 0 & 0 \\
 & \texttt{poll-duplex} & 3/3 & 25{,}600 & 0 & 4.8 & 10.5 & 13.9 & 100 & 37{,}505 & 870 & 8 & 263 & 256 & 51.4 & 0 & 0 \\
 & \texttt{poll-split} & 3/3 & 25{,}600 & 0 & 3.7 & 8.4 & 11.0 & 87 & 25{,}834 & 4{,}772 & 9 & 264 & 256 & 41.9 & 0 & 0 \\
 & \texttt{epoll} & 3/3 & 25{,}600 & 0 & 6.3 & 19.4 & 41.9 & 114 & 32{,}103 & 1{,}008 & 8 & 520 & 256 & 115.3 & 0 & 0 \\
 & \texttt{epoll-duplex} & 3/3 & 25{,}600 & 0 & 5.4 & 10.7 & 18.4 & 91 & 32{,}043 & 994 & 8 & 264 & 256 & 67.1 & 0 & 0 \\
 & \texttt{epoll-split} & 3/3 & 25{,}600 & 0 & 4.1 & 8.0 & 8.9 & 74 & 18{,}500 & 3{,}928 & 9 & 266 & 256 & 55.6 & 0 & 0 \\
\midrule
1{,}024 & \texttt{threads} & 3/3 & 15{,}114 & 1{,}745{,}710 & 18{,}253.6 & 18{,}643.5 & 18{,}643.5 & 193 & 88{,}807 & 2{,}938 & 1{,}032 & 2{,}054 & 1{,}024 & 293.6 & 0 & 0 \\
 & \texttt{pool-sender} & 3/3 & 10{,}370 & 1{,}840{,}600 & 19{,}327.4 & 19{,}503.3 & 19{,}503.3 & 192 & 91{,}182 & 2{,}820 & 1{,}033 & 2{,}055 & 1{,}024 & 310.4 & 0 & 0 \\
 & \texttt{peer-threads} & 3/3 & 20{,}794 & 1{,}632{,}110 & 16{,}374.6 & 17{,}551.8 & 17{,}551.8 & 190 & 106{,}962 & 3{,}968 & 1{,}032 & 2{,}054 & 1{,}024 & 159.4 & 0 & 0 \\
 & \texttt{poll} & 3/3 & 102{,}400 & 0 & 209.7 & 444.6 & 503.3 & 133 & 6{,}315 & 603 & 8 & 2{,}055 & 1{,}024 & 285.2 & 0 & 0 \\
 & \texttt{poll-duplex} & 3/3 & 102{,}400 & 0 & 75.5 & 167.8 & 247.5 & 145 & 15{,}858 & 745 & 8 & 1{,}031 & 1{,}024 & 94.4 & 0 & 0 \\
 & \texttt{poll-split} & 3/3 & 102{,}400 & 0 & 59.8 & 125.8 & 155.2 & 156 & 15{,}508 & 2{,}340 & 9 & 1{,}032 & 1{,}024 & 111.1 & 0 & 0 \\
 & \texttt{epoll} & 3/3 & 102{,}400 & 0 & 176.2 & 771.8 & 1{,}107.3 & 140 & 9{,}615 & 577 & 8 & 2{,}056 & 1{,}024 & 276.8 & 0 & 0 \\
 & \texttt{epoll-duplex} & 3/3 & 102{,}400 & 0 & 159.4 & 721.4 & 1{,}023.4 & 139 & 10{,}056 & 504 & 8 & 1{,}032 & 1{,}024 & 98.6 & 0 & 0 \\
 & \texttt{epoll-split} & 3/3 & 102{,}400 & 0 & 146.8 & 654.3 & 973.1 & 144 & 9{,}968 & 1{,}137 & 9 & 1{,}034 & 1{,}024 & 86.0 & 0 & 0 \\
\end{longtable}

\begin{longtable}{@{}llrrrrrrrrrrrrrrr@{}}
\caption{S2 (idle connections): all results.}\label{tab:full-S2}\\
\toprule
idle\_peers & Variant & valid & msg/s & lost & p50 & p99 & p99.9 & CPU \% & vol./s & invol./s & thr. & fds & ready & hs p50 & queued & evict. \\
\midrule
\endfirsthead
\toprule
idle\_peers & Variant & valid & msg/s & lost & p50 & p99 & p99.9 & CPU \% & vol./s & invol./s & thr. & fds & ready & hs p50 & queued & evict. \\
\midrule
\endhead
\bottomrule
\endfoot
0 & \texttt{threads} & 3/3 & 4{,}000 & 0 & 0.6 & 0.9 & 1.1 & 12 & 5{,}866 & 127 & 16 & 22 & 8 & 4.5 & 0 & 0 \\
 & \texttt{pool-sender} & 3/3 & 4{,}000 & 0 & 0.6 & 0.9 & 1.1 & 13 & 6{,}180 & 186 & 17 & 23 & 8 & 4.5 & 0 & 0 \\
 & \texttt{peer-threads} & 3/3 & 4{,}000 & 0 & 0.6 & 1.0 & 1.7 & 12 & 6{,}208 & 229 & 16 & 22 & 8 & 2.3 & 0 & 0 \\
 & \texttt{poll} & 3/3 & 4{,}000 & 0 & 0.7 & 1.0 & 1.3 & 19 & 4{,}314 & 8 & 8 & 23 & 8 & 3.7 & 0 & 0 \\
 & \texttt{poll-duplex} & 3/3 & 4{,}000 & 0 & 0.7 & 1.0 & 1.4 & 12 & 2{,}738 & 3 & 8 & 15 & 8 & 1.9 & 0 & 0 \\
 & \texttt{poll-split} & 3/3 & 4{,}000 & 0 & 0.7 & 0.9 & 1.3 & 10 & 2{,}103 & 26 & 9 & 16 & 8 & 1.6 & 0 & 0 \\
 & \texttt{epoll} & 3/3 & 4{,}000 & 0 & 0.6 & 1.0 & 1.2 & 17 & 3{,}936 & 6 & 8 & 24 & 8 & 4.5 & 0 & 0 \\
 & \texttt{epoll-duplex} & 3/3 & 4{,}000 & 0 & 0.7 & 1.0 & 1.6 & 13 & 2{,}899 & 2 & 8 & 16 & 8 & 2.4 & 0 & 0 \\
 & \texttt{epoll-split} & 3/3 & 4{,}000 & 0 & 0.7 & 1.0 & 1.4 & 11 & 2{,}421 & 13 & 9 & 18 & 8 & 2.1 & 0 & 0 \\
\midrule
100 & \texttt{threads} & 3/3 & 4{,}000 & 0 & 0.6 & 0.9 & 1.1 & 13 & 6{,}411 & 152 & 116 & 222 & 108 & 66.1 & 0 & 0 \\
 & \texttt{pool-sender} & 3/3 & 4{,}000 & 0 & 0.6 & 0.9 & 1.1 & 14 & 6{,}511 & 157 & 117 & 223 & 108 & 88.1 & 0 & 0 \\
 & \texttt{peer-threads} & 3/3 & 4{,}000 & 0 & 0.7 & 1.0 & 1.6 & 12 & 6{,}404 & 166 & 116 & 222 & 108 & 36.7 & 0 & 0 \\
 & \texttt{poll} & 3/3 & 4{,}000 & 0 & 0.7 & 1.2 & 1.5 & 35 & 6{,}962 & 78 & 8 & 223 & 108 & 81.8 & 0 & 0 \\
 & \texttt{poll-duplex} & 3/3 & 4{,}000 & 0 & 0.7 & 1.1 & 1.5 & 23 & 5{,}364 & 15 & 8 & 115 & 108 & 6.7 & 0 & 0 \\
 & \texttt{poll-split} & 3/3 & 4{,}000 & 0 & 0.7 & 1.0 & 1.4 & 20 & 5{,}230 & 415 & 9 & 116 & 108 & 8.3 & 0 & 0 \\
 & \texttt{epoll} & 3/3 & 4{,}000 & 0 & 0.6 & 1.0 & 1.3 & 18 & 3{,}963 & 5 & 8 & 224 & 108 & 71.3 & 0 & 0 \\
 & \texttt{epoll-duplex} & 3/3 & 4{,}000 & 0 & 0.7 & 1.0 & 1.8 & 13 & 2{,}909 & 2 & 8 & 116 & 108 & 8.3 & 0 & 0 \\
 & \texttt{epoll-split} & 3/3 & 4{,}000 & 0 & 0.7 & 1.0 & 1.5 & 10 & 1{,}999 & 9 & 9 & 118 & 108 & 8.4 & 0 & 0 \\
\midrule
300 & \texttt{threads} & 3/3 & 4{,}000 & 0 & 0.6 & 0.9 & 1.3 & 15 & 7{,}579 & 192 & 316 & 622 & 308 & 142.6 & 0 & 0 \\
 & \texttt{pool-sender} & 3/3 & 4{,}000 & 0 & 0.6 & 0.9 & 1.2 & 15 & 7{,}632 & 193 & 317 & 623 & 308 & 132.1 & 0 & 0 \\
 & \texttt{peer-threads} & 3/3 & 4{,}000 & 0 & 0.7 & 1.0 & 1.8 & 13 & 7{,}235 & 188 & 316 & 622 & 308 & 92.3 & 0 & 0 \\
 & \texttt{poll} & 3/3 & 4{,}000 & 0 & 0.8 & 1.4 & 2.2 & 51 & 8{,}096 & 11 & 8 & 623 & 308 & 142.6 & 0 & 0 \\
 & \texttt{poll-duplex} & 3/3 & 4{,}000 & 0 & 0.8 & 1.3 & 2.0 & 35 & 7{,}396 & 54 & 8 & 315 & 308 & 53.5 & 0 & 0 \\
 & \texttt{poll-split} & 3/3 & 4{,}000 & 0 & 0.8 & 1.1 & 1.8 & 34 & 6{,}609 & 116 & 9 & 316 & 308 & 62.9 & 0 & 0 \\
 & \texttt{epoll} & 3/3 & 4{,}000 & 0 & 0.6 & 1.0 & 1.3 & 18 & 3{,}966 & 8 & 8 & 624 & 308 & 134.2 & 0 & 0 \\
 & \texttt{epoll-duplex} & 3/3 & 4{,}000 & 0 & 0.6 & 1.0 & 1.7 & 13 & 2{,}888 & 3 & 8 & 316 & 308 & 55.6 & 0 & 0 \\
 & \texttt{epoll-split} & 3/3 & 4{,}000 & 0 & 0.6 & 1.0 & 1.6 & 10 & 2{,}044 & 15 & 9 & 318 & 308 & 54.5 & 0 & 0 \\
\midrule
1{,}000 & \texttt{threads} & 3/3 & 4{,}000 & 0 & 0.6 & 1.0 & 2.2 & 17 & 11{,}074 & 273 & 1{,}016 & 2{,}022 & 1{,}008 & 318.8 & 0 & 0 \\
 & \texttt{pool-sender} & 3/3 & 4{,}000 & 0 & 0.6 & 1.0 & 1.9 & 17 & 10{,}975 & 238 & 1{,}017 & 2{,}023 & 1{,}008 & 318.8 & 0 & 0 \\
 & \texttt{peer-threads} & 3/3 & 4{,}000 & 0 & 0.6 & 1.0 & 2.4 & 17 & 10{,}320 & 265 & 1{,}016 & 2{,}022 & 1{,}008 & 138.4 & 0 & 0 \\
 & \texttt{poll} & 3/3 & 4{,}000 & 0 & 1.4 & 3.3 & 4.8 & 100 & 7{,}215 & 12 & 8 & 2{,}023 & 1{,}008 & 268.4 & 0 & 0 \\
 & \texttt{poll-duplex} & 3/3 & 4{,}000 & 0 & 1.0 & 2.2 & 3.9 & 68 & 8{,}172 & 3 & 8 & 1{,}015 & 1{,}008 & 100.7 & 0 & 0 \\
 & \texttt{poll-split} & 3/3 & 4{,}000 & 0 & 0.9 & 2.2 & 3.4 & 65 & 8{,}114 & 9 & 9 & 1{,}016 & 1{,}008 & 90.2 & 0 & 0 \\
 & \texttt{epoll} & 3/3 & 4{,}000 & 0 & 0.6 & 1.0 & 2.7 & 17 & 4{,}028 & 6 & 8 & 2{,}024 & 1{,}008 & 276.8 & 0 & 0 \\
 & \texttt{epoll-duplex} & 3/3 & 4{,}000 & 0 & 0.7 & 1.1 & 2.8 & 13 & 2{,}906 & 2 & 8 & 1{,}016 & 1{,}008 & 83.9 & 0 & 0 \\
 & \texttt{epoll-split} & 3/3 & 4{,}000 & 0 & 0.6 & 1.0 & 1.9 & 10 & 1{,}987 & 8 & 9 & 1{,}018 & 1{,}008 & 86.0 & 0 & 0 \\
\end{longtable}

\begin{longtable}{@{}llrrrrrrrrrrrrrrr@{}}
\caption{S3 (slow consumers): all results.}\label{tab:full-S3}\\
\toprule
slow\_peers & Variant & valid & msg/s & lost & p50 & p99 & p99.9 & CPU \% & vol./s & invol./s & thr. & fds & ready & hs p50 & queued & evict. \\
\midrule
\endfirsthead
\toprule
slow\_peers & Variant & valid & msg/s & lost & p50 & p99 & p99.9 & CPU \% & vol./s & invol./s & thr. & fds & ready & hs p50 & queued & evict. \\
\midrule
\endhead
\bottomrule
\endfoot
0 & \texttt{threads} & 3/3 & 3{,}200 & 0 & 0.8 & 1.3 & 1.6 & 10 & 4{,}641 & 258 & 24 & 38 & 16 & 6.9 & 0 & 0 \\
 & \texttt{pool-sender} & 3/3 & 3{,}200 & 0 & 0.8 & 1.4 & 1.8 & 10 & 4{,}583 & 220 & 25 & 39 & 16 & 7.7 & 0 & 0 \\
 & \texttt{peer-threads} & 3/3 & 3{,}200 & 0 & 1.0 & 1.5 & 2.2 & 9 & 4{,}527 & 221 & 24 & 38 & 16 & 3.9 & 0 & 0 \\
 & \texttt{poll} & 3/3 & 3{,}200 & 0 & 1.0 & 1.5 & 2.0 & 14 & 2{,}919 & 11 & 8 & 39 & 16 & 8.1 & 0 & 0 \\
 & \texttt{poll-duplex} & 3/3 & 3{,}200 & 0 & 0.9 & 1.5 & 2.0 & 12 & 2{,}591 & 10 & 8 & 23 & 16 & 2.4 & 0 & 0 \\
 & \texttt{poll-split} & 3/3 & 3{,}200 & 0 & 0.9 & 1.4 & 2.2 & 9 & 1{,}422 & 70 & 9 & 24 & 16 & 2.8 & 0 & 0 \\
 & \texttt{epoll} & 3/3 & 3{,}200 & 0 & 0.9 & 1.5 & 2.0 & 13 & 2{,}638 & 12 & 8 & 40 & 16 & 6.8 & 0 & 0 \\
 & \texttt{epoll-duplex} & 3/3 & 3{,}200 & 0 & 1.0 & 1.5 & 2.0 & 10 & 2{,}351 & 8 & 8 & 24 & 16 & 2.5 & 0 & 0 \\
 & \texttt{epoll-split} & 3/3 & 3{,}200 & 0 & 0.9 & 1.4 & 2.0 & 8 & 1{,}393 & 43 & 9 & 26 & 16 & 2.7 & 0 & 0 \\
\midrule
1 & \texttt{threads} & 3/3 & 3{,}200 & 0 & 0.9 & 419.4 & 604.0 & 10 & 4{,}816 & 235 & 25 & 40 & 17 & 8.4 & 0 & 0 \\
 & \texttt{pool-sender} & 3/3 & 3{,}200 & 0 & 0.8 & 1.3 & 1.7 & 11 & 5{,}099 & 278 & 26 & 41 & 17 & 11.3 & 2{,}045 & 0 \\
 & \texttt{peer-threads} & 3/3 & 3{,}200 & 0 & 0.9 & 1.4 & 2.0 & 9 & 4{,}821 & 285 & 25 & 40 & 17 & 3.2 & 0 & 0 \\
 & \texttt{poll} & 3/3 & 3{,}200 & 0 & 1.0 & 1.5 & 1.9 & 16 & 3{,}294 & 17 & 8 & 41 & 17 & 8.1 & 2{,}046 & 0 \\
 & \texttt{poll-duplex} & 3/3 & 3{,}200 & 0 & 0.9 & 1.5 & 2.0 & 11 & 2{,}547 & 11 & 8 & 24 & 17 & 2.9 & 0 & 0 \\
 & \texttt{poll-split} & 3/3 & 3{,}200 & 0 & 0.9 & 1.4 & 2.6 & 8 & 1{,}413 & 68 & 9 & 25 & 17 & 2.2 & 0 & 0 \\
 & \texttt{epoll} & 3/3 & 3{,}200 & 0 & 0.9 & 1.5 & 1.8 & 13 & 2{,}700 & 10 & 8 & 42 & 17 & 11.0 & 2{,}046 & 0 \\
 & \texttt{epoll-duplex} & 3/3 & 3{,}200 & 0 & 0.9 & 1.5 & 2.0 & 11 & 2{,}277 & 6 & 8 & 25 & 17 & 2.9 & 0 & 0 \\
 & \texttt{epoll-split} & 3/3 & 3{,}200 & 0 & 0.9 & 1.5 & 1.9 & 8 & 1{,}301 & 34 & 9 & 27 & 17 & 3.9 & 0 & 0 \\
\midrule
4 & \texttt{threads} & 3/3 & 3{,}200 & 0 & 0.9 & 2{,}214.6 & 2{,}409.8 & 11 & 5{,}692 & 300 & 28 & 46 & 20 & 8.7 & 0 & 0 \\
 & \texttt{pool-sender} & 3/3 & 3{,}200 & 0 & 0.9 & 1.4 & 1.8 & 13 & 6{,}035 & 355 & 29 & 47 & 20 & 11.3 & 8{,}180 & 0 \\
 & \texttt{peer-threads} & 3/3 & 3{,}200 & 0 & 0.9 & 1.5 & 1.9 & 10 & 4{,}710 & 275 & 28 & 46 & 20 & 4.5 & 0 & 0 \\
 & \texttt{poll} & 3/3 & 3{,}200 & 0 & 1.0 & 1.6 & 1.9 & 19 & 3{,}918 & 31 & 8 & 47 & 20 & 8.9 & 8{,}184 & 0 \\
 & \texttt{poll-duplex} & 3/3 & 3{,}200 & 0 & 0.9 & 1.5 & 2.4 & 11 & 2{,}602 & 16 & 8 & 27 & 20 & 2.8 & 0 & 0 \\
 & \texttt{poll-split} & 3/3 & 3{,}200 & 0 & 0.9 & 1.5 & 2.2 & 8 & 1{,}344 & 85 & 9 & 28 & 20 & 3.3 & 0 & 0 \\
 & \texttt{epoll} & 3/3 & 3{,}200 & 0 & 1.0 & 1.6 & 1.9 & 16 & 3{,}427 & 68 & 8 & 48 & 20 & 9.4 & 8{,}178 & 0 \\
 & \texttt{epoll-duplex} & 3/3 & 3{,}200 & 0 & 1.0 & 1.5 & 1.9 & 11 & 2{,}379 & 9 & 8 & 28 & 20 & 2.6 & 0 & 0 \\
 & \texttt{epoll-split} & 3/3 & 3{,}200 & 0 & 0.9 & 1.5 & 2.1 & 7 & 1{,}224 & 30 & 9 & 30 & 20 & 3.1 & 0 & 0 \\
\end{longtable}

\begin{longtable}{@{}llrrrrrrrrrrrrrrr@{}}
\caption{S4 (churn): all results.}\label{tab:full-S4}\\
\toprule
churn & Variant & valid & msg/s & lost & p50 & p99 & p99.9 & CPU \% & vol./s & invol./s & thr. & fds & ready & hs p50 & queued & evict. \\
\midrule
\endfirsthead
\toprule
churn & Variant & valid & msg/s & lost & p50 & p99 & p99.9 & CPU \% & vol./s & invol./s & thr. & fds & ready & hs p50 & queued & evict. \\
\midrule
\endhead
\bottomrule
\endfoot
0 & \texttt{threads} & 3/3 & 3{,}200 & 0 & 2.2 & 4.0 & 5.4 & 10 & 4{,}104 & 290 & 72 & 134 & 64 & 55.6 & 0 & 0 \\
 & \texttt{pool-sender} & 3/3 & 3{,}200 & 0 & 2.3 & 3.9 & 4.8 & 10 & 4{,}107 & 290 & 73 & 135 & 64 & 53.5 & 0 & 0 \\
 & \texttt{peer-threads} & 3/3 & 3{,}200 & 0 & 2.3 & 4.0 & 5.5 & 9 & 4{,}333 & 268 & 72 & 134 & 64 & 24.6 & 0 & 0 \\
 & \texttt{poll} & 3/3 & 3{,}200 & 0 & 2.7 & 5.0 & 6.2 & 19 & 3{,}149 & 78 & 8 & 135 & 64 & 47.2 & 0 & 0 \\
 & \texttt{poll-duplex} & 3/3 & 3{,}200 & 0 & 2.3 & 4.2 & 5.2 & 12 & 2{,}111 & 52 & 8 & 71 & 64 & 4.8 & 0 & 0 \\
 & \texttt{poll-split} & 3/3 & 3{,}200 & 0 & 2.2 & 3.7 & 4.7 & 9 & 1{,}299 & 399 & 9 & 72 & 64 & 4.7 & 0 & 0 \\
 & \texttt{epoll} & 3/3 & 3{,}200 & 0 & 2.6 & 4.6 & 5.4 & 13 & 2{,}194 & 57 & 8 & 136 & 64 & 49.3 & 0 & 0 \\
 & \texttt{epoll-duplex} & 3/3 & 3{,}200 & 0 & 2.3 & 4.1 & 5.0 & 10 & 1{,}670 & 41 & 8 & 72 & 64 & 5.2 & 0 & 0 \\
 & \texttt{epoll-split} & 3/3 & 3{,}200 & 0 & 2.2 & 3.8 & 5.4 & 8 & 991 & 201 & 9 & 74 & 64 & 6.3 & 0 & 0 \\
\midrule
10 & \texttt{threads} & 3/3 & 3{,}190 & 200 & 2.2 & 4.3 & 5.4 & 11 & 4{,}305 & 363 & 72 & 135 & 64 & 3.8 & 0 & 0 \\
 & \texttt{pool-sender} & 3/3 & 3{,}189 & 226 & 2.4 & 67.1 & 104.9 & 12 & 4{,}381 & 379 & 73 & 136 & 64 & 4.2 & 0 & 0 \\
 & \texttt{peer-threads} & 3/3 & 3{,}189 & 222 & 2.3 & 64.0 & 104.9 & 10 & 4{,}531 & 356 & 72 & 135 & 64 & 3.0 & 0 & 0 \\
 & \texttt{poll} & 3/3 & 3{,}178 & 415 & 3.1 & 127.9 & 188.7 & 19 & 2{,}983 & 74 & 8 & 135 & 64 & 8.1 & 0 & 0 \\
 & \texttt{poll-duplex} & 3/3 & 3{,}190 & 204 & 2.4 & 4.7 & 83.9 & 13 & 2{,}295 & 73 & 8 & 71 & 64 & 3.1 & 0 & 0 \\
 & \texttt{poll-split} & 3/3 & 3{,}185 & 292 & 2.3 & 107.0 & 167.8 & 10 & 1{,}385 & 384 & 9 & 72 & 64 & 2.9 & 0 & 0 \\
 & \texttt{epoll} & 3/3 & 3{,}182 & 351 & 2.8 & 104.9 & 109.1 & 14 & 2{,}224 & 62 & 8 & 136 & 64 & 5.4 & 0 & 0 \\
 & \texttt{epoll-duplex} & 3/3 & 3{,}189 & 213 & 2.4 & 44.0 & 102.8 & 10 & 1{,}683 & 63 & 8 & 72 & 64 & 2.9 & 0 & 0 \\
 & \texttt{epoll-split} & 3/3 & 3{,}184 & 314 & 2.3 & 104.9 & 184.5 & 8 & 1{,}094 & 210 & 9 & 74 & 64 & 3.4 & 0 & 0 \\
\midrule
50 & \texttt{threads} & 3/3 & 3{,}150 & 1{,}000 & 2.5 & 4.6 & 5.6 & 14 & 4{,}641 & 359 & 72 & 135 & 64 & 3.5 & 0 & 0 \\
 & \texttt{pool-sender} & 3/3 & 3{,}129 & 1{,}412 & 2.8 & 104.9 & 107.0 & 16 & 4{,}873 & 355 & 73 & 136 & 64 & 3.7 & 0 & 0 \\
 & \texttt{peer-threads} & 3/3 & 3{,}133 & 1{,}293 & 2.4 & 102.8 & 146.8 & 14 & 4{,}873 & 384 & 72 & 135 & 64 & 2.6 & 0 & 0 \\
 & \texttt{poll} & 3/3 & 3{,}065 & 2{,}605 & 32.0 & 130.0 & 197.1 & 22 & 2{,}939 & 116 & 8 & 136 & 64 & 44.0 & 0 & 0 \\
 & \texttt{poll-duplex} & 3/3 & 3{,}143 & 1{,}140 & 2.6 & 86.0 & 125.8 & 14 & 2{,}377 & 154 & 8 & 71 & 64 & 2.8 & 0 & 0 \\
 & \texttt{poll-split} & 3/3 & 3{,}115 & 1{,}601 & 2.2 & 127.9 & 188.7 & 15 & 2{,}356 & 264 & 9 & 72 & 64 & 2.5 & 0 & 0 \\
 & \texttt{epoll} & 3/3 & 3{,}086 & 2{,}271 & 13.1 & 107.0 & 108.8 & 16 & 2{,}005 & 98 & 8 & 137 & 64 & 32.5 & 0 & 0 \\
 & \texttt{epoll-duplex} & 3/3 & 3{,}140 & 1{,}176 & 2.4 & 102.8 & 163.6 & 13 & 2{,}279 & 115 & 8 & 72 & 64 & 2.6 & 0 & 0 \\
 & \texttt{epoll-split} & 3/3 & 3{,}102 & 1{,}778 & 2.4 & 142.6 & 184.5 & 12 & 1{,}932 & 171 & 9 & 74 & 64 & 2.8 & 0 & 0 \\
\end{longtable}

\begin{longtable}{@{}llrrrrrrrrrrrrrrr@{}}
\caption{S5 (throughput): all results.}\label{tab:full-S5}\\
\toprule
size & Variant & valid & msg/s & lost & p50 & p99 & p99.9 & CPU \% & vol./s & invol./s & thr. & fds & ready & hs p50 & queued & evict. \\
\midrule
\endfirsthead
\toprule
size & Variant & valid & msg/s & lost & p50 & p99 & p99.9 & CPU \% & vol./s & invol./s & thr. & fds & ready & hs p50 & queued & evict. \\
\midrule
\endhead
\bottomrule
\endfoot
32 & \texttt{threads} & 3/3 & 186{,}244 & 0 & 0.6 & 1.5 & 2.3 & 105 & 42{,}966 & 7{,}578 & 16 & 22 & 8 & 4.6 & 0 & 0 \\
 & \texttt{pool-sender} & 3/3 & 183{,}801 & 0 & 0.6 & 1.5 & 2.5 & 101 & 42{,}044 & 7{,}245 & 17 & 23 & 8 & 6.4 & 0 & 0 \\
 & \texttt{peer-threads} & 3/3 & 203{,}039 & 0 & 0.6 & 1.3 & 2.0 & 112 & 53{,}287 & 9{,}252 & 16 & 22 & 8 & 2.9 & 0 & 0 \\
 & \texttt{poll} & 3/3 & 157{,}745 & 0 & 0.8 & 1.7 & 2.6 & 110 & 32{,}026 & 4{,}476 & 8 & 23 & 8 & 5.0 & 0 & 0 \\
 & \texttt{poll-duplex} & 3/3 & 187{,}580 & 0 & 0.6 & 1.4 & 2.0 & 111 & 34{,}182 & 5{,}795 & 8 & 15 & 8 & 2.4 & 0 & 0 \\
 & \texttt{poll-split} & 3/3 & 215{,}421 & 0 & 0.6 & 1.2 & 1.6 & 98 & 25{,}337 & 3{,}396 & 9 & 16 & 8 & 1.8 & 0 & 0 \\
 & \texttt{epoll} & 3/3 & 168{,}538 & 0 & 0.7 & 1.6 & 2.2 & 106 & 29{,}961 & 4{,}223 & 8 & 24 & 8 & 5.4 & 0 & 0 \\
 & \texttt{epoll-duplex} & 3/3 & 188{,}577 & 0 & 0.6 & 1.4 & 2.0 & 109 & 34{,}727 & 4{,}929 & 8 & 16 & 8 & 2.5 & 0 & 0 \\
 & \texttt{epoll-split} & 3/3 & 206{,}040 & 0 & 0.6 & 1.2 & 1.5 & 98 & 22{,}422 & 2{,}872 & 9 & 18 & 8 & 2.8 & 0 & 0 \\
\midrule
4{,}096 & \texttt{threads} & 3/3 & 113{,}314 & 0 & 1.0 & 2.1 & 2.8 & 127 & 29{,}282 & 6{,}993 & 16 & 22 & 8 & 4.8 & 0 & 0 \\
 & \texttt{pool-sender} & 3/3 & 109{,}301 & 0 & 1.1 & 2.1 & 2.8 & 128 & 29{,}797 & 6{,}540 & 17 & 23 & 8 & 4.0 & 0 & 0 \\
 & \texttt{peer-threads} & 3/3 & 104{,}707 & 0 & 1.1 & 2.4 & 3.1 & 137 & 31{,}793 & 7{,}916 & 16 & 22 & 8 & 2.7 & 0 & 0 \\
 & \texttt{poll} & 3/3 & 51{,}779 & 0 & 2.5 & 4.0 & 5.0 & 127 & 28{,}172 & 2{,}791 & 8 & 23 & 8 & 4.8 & 0 & 0 \\
 & \texttt{poll-duplex} & 3/3 & 56{,}075 & 0 & 2.3 & 3.7 & 4.7 & 126 & 28{,}286 & 3{,}184 & 8 & 15 & 8 & 2.3 & 0 & 0 \\
 & \texttt{poll-split} & 3/3 & 82{,}008 & 0 & 1.4 & 3.4 & 4.5 & 145 & 25{,}535 & 10{,}037 & 9 & 16 & 8 & 2.8 & 0 & 0 \\
 & \texttt{epoll} & 3/3 & 54{,}707 & 0 & 2.3 & 4.3 & 5.2 & 128 & 26{,}282 & 3{,}212 & 8 & 24 & 8 & 4.5 & 0 & 0 \\
 & \texttt{epoll-duplex} & 3/3 & 51{,}111 & 0 & 2.5 & 4.2 & 4.8 & 126 & 29{,}110 & 3{,}175 & 8 & 16 & 8 & 2.9 & 0 & 0 \\
 & \texttt{epoll-split} & 3/3 & 90{,}040 & 0 & 1.3 & 2.6 & 3.5 & 144 & 26{,}010 & 10{,}226 & 9 & 18 & 8 & 2.2 & 0 & 0 \\
\midrule
65{,}536 & \texttt{threads} & 3/3 & 12{,}405 & 0 & 10.0 & 13.9 & 19.9 & 173 & 13{,}146 & 5{,}238 & 16 & 22 & 8 & 8.1 & 0 & 0 \\
 & \texttt{pool-sender} & 3/3 & 17{,}687 & 0 & 4.5 & 9.7 & 12.1 & 170 & 16{,}557 & 4{,}925 & 17 & 23 & 8 & 4.7 & 52 & 0 \\
 & \texttt{peer-threads} & 3/3 & 16{,}152 & 0 & 6.6 & 12.1 & 15.2 & 179 & 12{,}752 & 4{,}429 & 16 & 22 & 8 & 2.9 & 31 & 0 \\
 & \texttt{poll} & 3/3 & 7{,}967 & 0 & 16.3 & 18.9 & 22.5 & 122 & 16{,}751 & 378 & 8 & 23 & 8 & 5.1 & 34 & 0 \\
 & \texttt{poll-duplex} & 3/3 & 8{,}590 & 0 & 14.9 & 17.3 & 21.0 & 124 & 17{,}093 & 453 & 8 & 15 & 8 & 2.3 & 35 & 0 \\
 & \texttt{poll-split} & 3/3 & 10{,}388 & 0 & 12.3 & 15.2 & 20.4 & 140 & 14{,}292 & 932 & 9 & 16 & 8 & 2.2 & 26 & 0 \\
 & \texttt{epoll} & 3/3 & 8{,}457 & 0 & 15.2 & 17.8 & 23.6 & 123 & 15{,}774 & 508 & 8 & 24 & 8 & 3.1 & 50 & 0 \\
 & \texttt{epoll-duplex} & 3/3 & 7{,}627 & 0 & 16.8 & 19.4 & 25.7 & 144 & 17{,}102 & 1{,}519 & 8 & 16 & 8 & 2.4 & 5 & 0 \\
 & \texttt{epoll-split} & 3/3 & 10{,}658 & 0 & 12.1 & 14.9 & 20.4 & 138 & 12{,}392 & 821 & 9 & 18 & 8 & 2.0 & 15 & 0 \\
\end{longtable}

\begin{longtable}{@{}llrrrrrrrrrrrrrrr@{}}
\caption{S6a (WAN churn): all results.}\label{tab:full-S6a}\\
\toprule
rtt\_ms & Variant & valid & msg/s & lost & p50 & p99 & p99.9 & CPU \% & vol./s & invol./s & thr. & fds & ready & hs p50 & queued & evict. \\
\midrule
\endfirsthead
\toprule
rtt\_ms & Variant & valid & msg/s & lost & p50 & p99 & p99.9 & CPU \% & vol./s & invol./s & thr. & fds & ready & hs p50 & queued & evict. \\
\midrule
\endhead
\bottomrule
\endfoot
0 & \texttt{threads} & 3/3 & 3{,}196 & 80 & 0.7 & 1.5 & 2.8 & 10 & 4{,}394 & 384 & 40 & 70 & 32 & 0.6 & 0 & 0 \\
 & \texttt{pool-sender} & 3/3 & 3{,}196 & 80 & 0.7 & 1.6 & 2.8 & 10 & 4{,}483 & 403 & 41 & 71 & 32 & 0.6 & 0 & 0 \\
 & \texttt{peer-threads} & 3/3 & 3{,}192 & 154 & 0.7 & 83.9 & 102.0 & 10 & 4{,}514 & 363 & 40 & 71 & 32 & 0.5 & 0 & 0 \\
 & \texttt{poll} & 3/3 & 3{,}186 & 274 & 1.1 & 92.3 & 102.6 & 19 & 3{,}423 & 88 & 8 & 71 & 32 & 0.7 & 0 & 0 \\
 & \texttt{poll-duplex} & 3/3 & 3{,}194 & 120 & 0.7 & 73.4 & 101.4 & 13 & 2{,}285 & 30 & 8 & 39 & 32 & 0.5 & 0 & 0 \\
 & \texttt{poll-split} & 3/3 & 3{,}189 & 224 & 0.5 & 92.3 & 102.8 & 7 & 764 & 88 & 9 & 40 & 32 & 0.5 & 0 & 0 \\
 & \texttt{epoll} & 3/3 & 3{,}192 & 166 & 0.8 & 83.9 & 102.6 & 15 & 2{,}700 & 78 & 8 & 72 & 32 & 0.6 & 0 & 0 \\
 & \texttt{epoll-duplex} & 3/3 & 3{,}194 & 121 & 0.6 & 73.4 & 101.3 & 10 & 1{,}560 & 29 & 8 & 40 & 32 & 0.5 & 0 & 0 \\
 & \texttt{epoll-split} & 3/3 & 3{,}190 & 196 & 0.5 & 92.3 & 102.4 & 7 & 708 & 59 & 9 & 42 & 32 & 0.5 & 0 & 0 \\
\midrule
4 & \texttt{threads} & 3/3 & 3{,}190 & 195 & 4.7 & 9.2 & 9.4 & 8 & 4{,}325 & 232 & 40 & 70 & 32 & 13.1 & 0 & 0 \\
 & \texttt{pool-sender} & 3/3 & 3{,}189 & 205 & 4.7 & 9.2 & 98.6 & 8 & 4{,}406 & 213 & 41 & 71 & 32 & 13.1 & 0 & 0 \\
 & \texttt{peer-threads} & 3/3 & 3{,}189 & 226 & 4.7 & 75.5 & 96.5 & 8 & 4{,}403 & 255 & 40 & 70 & 32 & 8.9 & 0 & 0 \\
 & \texttt{poll} & 3/3 & 3{,}171 & 307 & 4.5 & 107.0 & 111.1 & 17 & 4{,}312 & 92 & 8 & 71 & 32 & 13.1 & 0 & 0 \\
 & \texttt{poll-duplex} & 3/3 & 3{,}190 & 208 & 4.7 & 5.8 & 75.5 & 10 & 2{,}838 & 124 & 8 & 39 & 32 & 8.7 & 0 & 0 \\
 & \texttt{poll-split} & 3/3 & 3{,}188 & 244 & 4.6 & 86.0 & 104.9 & 7 & 1{,}447 & 392 & 9 & 40 & 32 & 8.7 & 0 & 0 \\
 & \texttt{epoll} & 3/3 & 3{,}180 & 257 & 4.5 & 104.9 & 111.1 & 12 & 3{,}234 & 87 & 8 & 72 & 32 & 13.1 & 0 & 0 \\
 & \texttt{epoll-duplex} & 3/3 & 3{,}190 & 205 & 4.7 & 5.8 & 88.1 & 9 & 2{,}677 & 141 & 8 & 40 & 32 & 8.7 & 0 & 0 \\
 & \texttt{epoll-split} & 3/3 & 3{,}186 & 274 & 4.6 & 96.5 & 106.1 & 6 & 1{,}203 & 262 & 9 & 42 & 32 & 8.7 & 0 & 0 \\
\midrule
20 & \texttt{threads} & 3/3 & 3{,}158 & 792 & 21.5 & 41.9 & 43.0 & 9 & 4{,}172 & 179 & 40 & 70 & 32 & 61.9 & 0 & 0 \\
 & \texttt{pool-sender} & 3/3 & 3{,}156 & 856 & 21.5 & 75.5 & 163.6 & 10 & 4{,}419 & 202 & 41 & 71 & 32 & 61.9 & 0 & 0 \\
 & \texttt{peer-threads} & 3/3 & 3{,}161 & 735 & 21.0 & 83.9 & 121.3 & 10 & 4{,}890 & 246 & 40 & 70 & 32 & 40.9 & 0 & 0 \\
 & \texttt{poll} & 3/3 & 3{,}099 & 1{,}305 & 21.0 & 172.0 & 180.4 & 18 & 3{,}814 & 27 & 8 & 71 & 32 & 61.9 & 0 & 0 \\
 & \texttt{poll-duplex} & 3/3 & 3{,}158 & 750 & 21.0 & 102.8 & 121.2 & 12 & 2{,}817 & 50 & 8 & 39 & 32 & 40.9 & 0 & 0 \\
 & \texttt{poll-split} & 3/3 & 3{,}156 & 749 & 20.4 & 102.8 & 121.5 & 11 & 2{,}237 & 196 & 9 & 40 & 32 & 40.9 & 0 & 0 \\
 & \texttt{epoll} & 3/3 & 3{,}110 & 1{,}235 & 21.0 & 172.0 & 180.4 & 12 & 2{,}680 & 17 & 8 & 72 & 32 & 61.9 & 0 & 0 \\
 & \texttt{epoll-duplex} & 3/3 & 3{,}162 & 727 & 21.0 & 52.4 & 121.2 & 9 & 2{,}225 & 32 & 8 & 40 & 32 & 40.9 & 0 & 0 \\
 & \texttt{epoll-split} & 3/3 & 3{,}156 & 750 & 21.0 & 102.8 & 121.1 & 8 & 1{,}558 & 76 & 9 & 42 & 32 & 40.9 & 0 & 0 \\
\midrule
30 & \texttt{threads} & 3/3 & 3{,}147 & 1{,}060 & 34.6 & 61.9 & 75.5 & 8 & 4{,}034 & 110 & 40 & 70 & 32 & 92.3 & 0 & 0 \\
 & \texttt{pool-sender} & 3/3 & 3{,}147 & 1{,}064 & 34.6 & 64.0 & 83.9 & 10 & 4{,}531 & 119 & 41 & 71 & 32 & 92.3 & 0 & 0 \\
 & \texttt{peer-threads} & 3/3 & 3{,}141 & 1{,}115 & 30.9 & 90.2 & 130.0 & 10 & 4{,}880 & 244 & 40 & 70 & 32 & 60.8 & 0 & 0 \\
 & \texttt{poll} & 3/3 & 2{,}801 & 3{,}960 & 107.0 & 251.7 & 276.8 & 15 & 3{,}069 & 12 & 8 & 71 & 32 & 163.6 & 0 & 0 \\
 & \texttt{poll-duplex} & 3/3 & 3{,}135 & 1{,}174 & 30.9 & 113.2 & 131.8 & 12 & 3{,}458 & 85 & 8 & 39 & 32 & 60.8 & 0 & 0 \\
 & \texttt{poll-split} & 3/3 & 3{,}127 & 1{,}249 & 30.4 & 121.6 & 132.1 & 10 & 2{,}414 & 236 & 9 & 40 & 32 & 60.8 & 0 & 0 \\
 & \texttt{epoll} & 3/3 & 2{,}845 & 3{,}618 & 88.1 & 251.7 & 275.7 & 11 & 2{,}444 & 11 & 8 & 72 & 32 & 155.2 & 0 & 0 \\
 & \texttt{epoll-duplex} & 3/3 & 3{,}137 & 1{,}163 & 30.9 & 104.9 & 131.1 & 10 & 2{,}793 & 26 & 8 & 40 & 32 & 60.8 & 0 & 0 \\
 & \texttt{epoll-split} & 3/3 & 3{,}135 & 1{,}181 & 30.4 & 111.1 & 131.9 & 8 & 1{,}808 & 119 & 9 & 42 & 32 & 60.8 & 0 & 0 \\
\midrule
40 & \texttt{threads} & 0/3 & -- & -- & -- & -- & -- & -- & -- & -- & -- & -- & -- & -- & -- & -- \\
 & \texttt{pool-sender} & 0/3 & -- & -- & -- & -- & -- & -- & -- & -- & -- & -- & -- & -- & -- & -- \\
 & \texttt{peer-threads} & 3/3 & 3{,}129 & 1{,}385 & 40.9 & 56.6 & 132.1 & 10 & 4{,}727 & 264 & 40 & 70 & 32 & 81.8 & 0 & 0 \\
 & \texttt{poll} & 0/3 & -- & -- & -- & -- & -- & -- & -- & -- & -- & -- & -- & -- & -- & -- \\
 & \texttt{poll-duplex} & 3/3 & 3{,}124 & 1{,}428 & 40.9 & 113.2 & 142.1 & 11 & 2{,}641 & 40 & 8 & 39 & 32 & 81.8 & 0 & 0 \\
 & \texttt{poll-split} & 3/3 & 3{,}116 & 1{,}501 & 40.9 & 130.0 & 142.2 & 11 & 2{,}189 & 174 & 9 & 40 & 32 & 81.8 & 0 & 0 \\
 & \texttt{epoll} & 0/3 & -- & -- & -- & -- & -- & -- & -- & -- & -- & -- & -- & -- & -- & -- \\
 & \texttt{epoll-duplex} & 3/3 & 3{,}121 & 1{,}448 & 40.9 & 123.7 & 141.6 & 10 & 2{,}597 & 29 & 8 & 40 & 32 & 81.8 & 0 & 0 \\
 & \texttt{epoll-split} & 3/3 & 3{,}121 & 1{,}441 & 40.9 & 123.7 & 142.0 & 9 & 1{,}854 & 77 & 9 & 42 & 32 & 81.8 & 0 & 0 \\
\midrule
50 & \texttt{threads} & 0/3 & -- & -- & -- & -- & -- & -- & -- & -- & -- & -- & -- & -- & -- & -- \\
 & \texttt{pool-sender} & 0/3 & -- & -- & -- & -- & -- & -- & -- & -- & -- & -- & -- & -- & -- & -- \\
 & \texttt{peer-threads} & 3/3 & 3{,}103 & 1{,}828 & 51.4 & 115.3 & 151.0 & 10 & 4{,}957 & 237 & 40 & 70 & 32 & 102.8 & 0 & 0 \\
 & \texttt{poll} & 0/3 & -- & -- & -- & -- & -- & -- & -- & -- & -- & -- & -- & -- & -- & -- \\
 & \texttt{poll-duplex} & 3/3 & 3{,}103 & 1{,}823 & 51.4 & 119.5 & 151.0 & 12 & 3{,}336 & 54 & 8 & 39 & 32 & 102.8 & 0 & 0 \\
 & \texttt{poll-split} & 3/3 & 3{,}097 & 1{,}871 & 51.4 & 142.6 & 151.0 & 10 & 2{,}255 & 255 & 9 & 40 & 32 & 100.7 & 0 & 0 \\
 & \texttt{epoll} & 0/3 & -- & -- & -- & -- & -- & -- & -- & -- & -- & -- & -- & -- & -- & -- \\
 & \texttt{epoll-duplex} & 3/3 & 3{,}107 & 1{,}802 & 51.4 & 102.8 & 151.0 & 10 & 2{,}591 & 37 & 8 & 40 & 32 & 100.7 & 0 & 0 \\
 & \texttt{epoll-split} & 3/3 & 3{,}100 & 1{,}833 & 51.4 & 132.1 & 151.0 & 8 & 1{,}710 & 101 & 9 & 42 & 32 & 100.7 & 0 & 0 \\
\end{longtable}

\begin{longtable}{@{}llrrrrrrrrrrrrrrr@{}}
\caption{S6c (WAN churn, halved): all results.}\label{tab:full-S6c}\\
\toprule
rtt\_ms & Variant & valid & msg/s & lost & p50 & p99 & p99.9 & CPU \% & vol./s & invol./s & thr. & fds & ready & hs p50 & queued & evict. \\
\midrule
\endfirsthead
\toprule
rtt\_ms & Variant & valid & msg/s & lost & p50 & p99 & p99.9 & CPU \% & vol./s & invol./s & thr. & fds & ready & hs p50 & queued & evict. \\
\midrule
\endhead
\bottomrule
\endfoot
50 & \texttt{threads} & 3/3 & 3{,}022 & 1{,}196 & 51.4 & 100.7 & 101.5 & 11 & 5{,}115 & 150 & 40 & 70 & 32 & 155.2 & 0 & 0 \\
 & \texttt{pool-sender} & 3/3 & 3{,}022 & 1{,}205 & 51.4 & 100.7 & 102.1 & 11 & 5{,}191 & 161 & 41 & 71 & 32 & 155.2 & 0 & 0 \\
 & \texttt{peer-threads} & 3/3 & 3{,}148 & 939 & 51.4 & 121.6 & 151.0 & 10 & 4{,}823 & 238 & 40 & 70 & 32 & 102.8 & 0 & 0 \\
 & \texttt{poll} & 3/3 & 3{,}018 & 1{,}237 & 51.4 & 192.9 & 285.2 & 16 & 3{,}655 & 34 & 8 & 71 & 32 & 155.2 & 0 & 0 \\
 & \texttt{poll-duplex} & 3/3 & 3{,}149 & 932 & 51.4 & 119.5 & 151.0 & 11 & 3{,}130 & 47 & 8 & 39 & 32 & 102.8 & 0 & 0 \\
 & \texttt{poll-split} & 3/3 & 3{,}142 & 944 & 51.4 & 132.1 & 151.0 & 9 & 2{,}098 & 284 & 9 & 40 & 32 & 102.8 & 0 & 0 \\
 & \texttt{epoll} & 3/3 & 3{,}014 & 1{,}293 & 51.4 & 226.5 & 293.6 & 11 & 2{,}602 & 7 & 8 & 72 & 32 & 155.2 & 0 & 0 \\
 & \texttt{epoll-duplex} & 3/3 & 3{,}149 & 951 & 51.4 & 102.8 & 151.0 & 9 & 2{,}697 & 37 & 8 & 40 & 32 & 102.8 & 0 & 0 \\
 & \texttt{epoll-split} & 3/3 & 3{,}145 & 956 & 51.4 & 132.1 & 151.0 & 7 & 1{,}515 & 121 & 9 & 42 & 32 & 100.7 & 0 & 0 \\
\end{longtable}

\begin{longtable}{@{}llrrrrrrrrrrrrrrr@{}}
\caption{S6b (WAN bottleneck): all results.}\label{tab:full-S6b}\\
\toprule
rate\_mbit & Variant & valid & msg/s & lost & p50 & p99 & p99.9 & CPU \% & vol./s & invol./s & thr. & fds & ready & hs p50 & queued & evict. \\
\midrule
\endfirsthead
\toprule
rate\_mbit & Variant & valid & msg/s & lost & p50 & p99 & p99.9 & CPU \% & vol./s & invol./s & thr. & fds & ready & hs p50 & queued & evict. \\
\midrule
\endhead
\bottomrule
\endfoot
50 & \texttt{threads} & 3/3 & 352 & 0 & 671.1 & 1{,}778.4 & 2{,}617.2 & 4 & 840 & 0 & 40 & 70 & 32 & 77.6 & 0 & 0 \\
 & \texttt{pool-sender} & 3/3 & 361 & 0 & 511.7 & 2{,}818.6 & 3{,}758.1 & 6 & 1{,}368 & 0 & 41 & 71 & 32 & 77.6 & 4{,}775 & 0 \\
 & \texttt{peer-threads} & 3/3 & 363 & 0 & 637.5 & 1{,}744.8 & 2{,}415.9 & 6 & 1{,}490 & 0 & 40 & 70 & 32 & 10.0 & 5{,}431 & 0 \\
 & \texttt{poll} & 3/3 & 361 & 0 & 528.5 & 2{,}818.6 & 4{,}160.8 & 8 & 1{,}324 & 0 & 8 & 71 & 32 & 77.6 & 4{,}544 & 0 \\
 & \texttt{poll-duplex} & 3/3 & 363 & 0 & 637.5 & 1{,}744.8 & 2{,}348.8 & 7 & 1{,}341 & 0 & 8 & 39 & 32 & 9.2 & 5{,}306 & 0 \\
 & \texttt{poll-split} & 3/3 & 363 & 0 & 637.5 & 1{,}811.9 & 2{,}415.9 & 7 & 1{,}206 & 0 & 9 & 40 & 32 & 9.4 & 5{,}360 & 0 \\
 & \texttt{epoll} & 3/3 & 361 & 0 & 520.1 & 2{,}885.7 & 3{,}959.4 & 6 & 1{,}298 & 0 & 8 & 72 & 32 & 77.6 & 4{,}253 & 0 \\
 & \texttt{epoll-duplex} & 3/3 & 363 & 0 & 637.5 & 1{,}845.5 & 2{,}684.3 & 6 & 1{,}349 & 0 & 8 & 40 & 32 & 9.4 & 5{,}390 & 0 \\
 & \texttt{epoll-split} & 3/3 & 363 & 0 & 637.5 & 1{,}845.5 & 2{,}348.8 & 5 & 1{,}210 & 0 & 9 & 42 & 32 & 9.4 & 5{,}281 & 0 \\
\midrule
200 & \texttt{threads} & 3/3 & 1{,}353 & 0 & 163.6 & 637.5 & 1{,}174.4 & 10 & 2{,}810 & 0 & 40 & 70 & 32 & 75.5 & 0 & 0 \\
 & \texttt{pool-sender} & 3/3 & 1{,}446 & 0 & 138.4 & 922.7 & 1{,}778.4 & 13 & 4{,}395 & 1 & 41 & 71 & 32 & 77.6 & 2{,}041 & 0 \\
 & \texttt{peer-threads} & 3/3 & 1{,}452 & 0 & 151.0 & 637.5 & 1{,}057.0 & 14 & 4{,}507 & 1 & 40 & 70 & 32 & 10.5 & 2{,}331 & 0 \\
 & \texttt{poll} & 3/3 & 1{,}446 & 0 & 146.8 & 872.4 & 1{,}577.1 & 18 & 4{,}347 & 0 & 8 & 71 & 32 & 77.6 & 1{,}794 & 0 \\
 & \texttt{poll-duplex} & 3/3 & 1{,}452 & 0 & 151.0 & 620.8 & 1{,}107.3 & 16 & 4{,}353 & 0 & 8 & 39 & 32 & 9.2 & 1{,}895 & 0 \\
 & \texttt{poll-split} & 3/3 & 1{,}452 & 0 & 151.0 & 637.5 & 1{,}023.4 & 16 & 4{,}329 & 1 & 9 & 40 & 32 & 8.9 & 2{,}072 & 0 \\
 & \texttt{epoll} & 3/3 & 1{,}446 & 0 & 138.4 & 872.4 & 1{,}610.6 & 13 & 4{,}360 & 0 & 8 & 72 & 32 & 77.6 & 2{,}716 & 0 \\
 & \texttt{epoll-duplex} & 3/3 & 1{,}452 & 0 & 151.0 & 604.0 & 1{,}057.0 & 13 & 4{,}355 & 0 & 8 & 40 & 32 & 9.4 & 2{,}071 & 0 \\
 & \texttt{epoll-split} & 3/3 & 1{,}451 & 0 & 151.0 & 620.8 & 1{,}140.9 & 13 & 4{,}339 & 0 & 9 & 42 & 32 & 8.9 & 1{,}963 & 0 \\
\midrule
1{,}000 & \texttt{threads} & 3/3 & 7{,}220 & 0 & 33.6 & 121.6 & 192.9 & 44 & 19{,}465 & 47 & 40 & 70 & 32 & 77.6 & 0 & 0 \\
 & \texttt{pool-sender} & 3/3 & 7{,}221 & 0 & 33.6 & 119.5 & 188.7 & 44 & 19{,}338 & 58 & 41 & 71 & 32 & 77.6 & 23 & 0 \\
 & \texttt{peer-threads} & 3/3 & 7{,}252 & 0 & 36.7 & 96.5 & 130.0 & 48 & 19{,}417 & 79 & 40 & 70 & 32 & 10.2 & 21 & 0 \\
 & \texttt{poll} & 3/3 & 7{,}224 & 0 & 34.6 & 119.5 & 180.4 & 59 & 18{,}652 & 22 & 8 & 71 & 32 & 77.6 & 9 & 0 \\
 & \texttt{poll-duplex} & 3/3 & 7{,}252 & 0 & 33.6 & 94.4 & 138.4 & 55 & 18{,}107 & 31 & 8 & 39 & 32 & 9.4 & 19 & 0 \\
 & \texttt{poll-split} & 3/3 & 7{,}252 & 0 & 35.7 & 98.6 & 151.0 & 54 & 17{,}384 & 45 & 9 & 40 & 32 & 9.1 & 7 & 0 \\
 & \texttt{epoll} & 3/3 & 7{,}222 & 0 & 33.6 & 117.4 & 192.9 & 45 & 18{,}255 & 25 & 8 & 72 & 32 & 77.6 & 51 & 0 \\
 & \texttt{epoll-duplex} & 3/3 & 7{,}251 & 0 & 36.7 & 98.6 & 142.6 & 46 & 18{,}068 & 26 & 8 & 40 & 32 & 9.2 & 9 & 0 \\
 & \texttt{epoll-split} & 3/3 & 7{,}250 & 0 & 36.7 & 98.6 & 142.6 & 46 & 17{,}679 & 22 & 9 & 42 & 32 & 8.9 & 17 & 0 \\
\end{longtable}

\begin{longtable}{@{}llrrrrrrrrrrrrrrr@{}}
\caption{S7 (capped fan-in): all results.}\label{tab:full-S7}\\
\toprule
max\_peers & Variant & valid & msg/s & lost & p50 & p99 & p99.9 & CPU \% & vol./s & invol./s & thr. & fds & ready & hs p50 & queued & evict. \\
\midrule
\endfirsthead
\toprule
max\_peers & Variant & valid & msg/s & lost & p50 & p99 & p99.9 & CPU \% & vol./s & invol./s & thr. & fds & ready & hs p50 & queued & evict. \\
\midrule
\endhead
\bottomrule
\endfoot
64 & \texttt{threads} & 3/3 & 6{,}400 & 0 & 1.8 & 3.3 & 4.5 & 18 & 8{,}189 & 669 & 72 & 134 & 64 & 343.9 & 0 & 960 \\
 & \texttt{pool-sender} & 3/3 & 6{,}400 & 0 & 1.8 & 3.0 & 4.1 & 17 & 8{,}432 & 705 & 73 & 135 & 64 & 310.4 & 0 & 960 \\
 & \texttt{peer-threads} & 3/3 & 6{,}400 & 0 & 1.8 & 3.1 & 4.3 & 16 & 8{,}499 & 697 & 72 & 134 & 64 & 107.0 & 0 & 960 \\
 & \texttt{poll} & 3/3 & 6{,}400 & 0 & 2.6 & 5.0 & 5.6 & 45 & 8{,}108 & 171 & 8 & 135 & 64 & 172.0 & 0 & 0 \\
 & \texttt{poll-duplex} & 3/3 & 8{,}000 & 0 & 2.4 & 4.3 & 5.1 & 33 & 7{,}937 & 228 & 8 & 87 & 80 & 55.6 & 0 & 0 \\
 & \texttt{poll-split} & 3/3 & 8{,}000 & 0 & 2.0 & 2.9 & 3.8 & 23 & 4{,}780 & 1{,}669 & 9 & 88 & 80 & 58.7 & 0 & 0 \\
 & \texttt{epoll} & 3/3 & 6{,}300 & 0 & 2.3 & 3.9 & 4.7 & 27 & 4{,}980 & 112 & 8 & 134 & 63 & 151.0 & 0 & 0 \\
 & \texttt{epoll-duplex} & 3/3 & 8{,}000 & 0 & 2.2 & 4.1 & 4.7 & 25 & 5{,}867 & 198 & 8 & 88 & 80 & 44.0 & 0 & 0 \\
 & \texttt{epoll-split} & 3/3 & 8{,}000 & 0 & 1.9 & 2.9 & 3.9 & 17 & 3{,}230 & 1{,}052 & 9 & 90 & 80 & 73.4 & 0 & 0 \\
\midrule
128 & \texttt{threads} & 3/3 & 12{,}800 & 0 & 2.9 & 5.4 & 7.6 & 34 & 15{,}401 & 1{,}311 & 136 & 262 & 128 & 377.5 & 0 & 896 \\
 & \texttt{pool-sender} & 3/3 & 12{,}800 & 0 & 2.9 & 5.4 & 7.6 & 34 & 15{,}492 & 1{,}387 & 137 & 263 & 128 & 335.5 & 0 & 896 \\
 & \texttt{peer-threads} & 3/3 & 12{,}800 & 0 & 2.9 & 4.6 & 5.9 & 32 & 16{,}682 & 1{,}767 & 136 & 262 & 128 & 119.5 & 0 & 896 \\
 & \texttt{poll} & 3/3 & 12{,}700 & 0 & 3.9 & 7.7 & 9.7 & 78 & 19{,}519 & 329 & 8 & 261 & 127 & 218.1 & 0 & 0 \\
 & \texttt{poll-duplex} & 3/3 & 14{,}400 & 0 & 3.4 & 6.6 & 7.5 & 60 & 16{,}504 & 486 & 8 & 151 & 144 & 52.4 & 0 & 0 \\
 & \texttt{poll-split} & 3/3 & 14{,}400 & 0 & 2.9 & 5.4 & 6.2 & 51 & 12{,}996 & 3{,}193 & 9 & 152 & 144 & 77.6 & 0 & 0 \\
 & \texttt{epoll} & 3/3 & 12{,}800 & 0 & 3.5 & 6.4 & 7.3 & 56 & 12{,}247 & 454 & 8 & 264 & 128 & 247.5 & 0 & 0 \\
 & \texttt{epoll-duplex} & 3/3 & 14{,}400 & 0 & 3.5 & 6.6 & 7.5 & 52 & 13{,}348 & 494 & 8 & 152 & 144 & 67.1 & 0 & 0 \\
 & \texttt{epoll-split} & 3/3 & 14{,}400 & 0 & 2.8 & 5.0 & 6.7 & 36 & 7{,}371 & 2{,}162 & 9 & 154 & 144 & 115.3 & 0 & 0 \\
\midrule
256 & \texttt{threads} & 3/3 & 25{,}600 & 0 & 4.1 & 7.1 & 9.4 & 50 & 28{,}481 & 1{,}099 & 264 & 518 & 256 & 360.7 & 0 & 768 \\
 & \texttt{pool-sender} & 3/3 & 25{,}600 & 0 & 4.1 & 7.2 & 9.2 & 51 & 28{,}284 & 1{,}151 & 265 & 519 & 256 & 335.5 & 0 & 768 \\
 & \texttt{peer-threads} & 3/3 & 25{,}600 & 0 & 4.1 & 6.8 & 10.7 & 58 & 30{,}095 & 1{,}949 & 264 & 518 & 256 & 146.8 & 0 & 768 \\
 & \texttt{poll} & 3/3 & 25{,}600 & 0 & 9.4 & 19.9 & 29.9 & 131 & 28{,}249 & 398 & 8 & 519 & 256 & 335.5 & 0 & 0 \\
 & \texttt{poll-duplex} & 3/3 & 27{,}200 & 0 & 5.2 & 11.8 & 15.7 & 115 & 40{,}956 & 929 & 8 & 279 & 272 & 90.2 & 0 & 0 \\
 & \texttt{poll-split} & 3/3 & 27{,}200 & 0 & 3.9 & 9.2 & 11.8 & 93 & 27{,}941 & 5{,}033 & 9 & 280 & 272 & 104.9 & 0 & 0 \\
 & \texttt{epoll} & 3/3 & 25{,}600 & 0 & 6.4 & 22.0 & 38.8 & 115 & 32{,}561 & 1{,}050 & 8 & 520 & 256 & 310.4 & 0 & 0 \\
 & \texttt{epoll-duplex} & 3/3 & 27{,}200 & 0 & 5.5 & 11.8 & 22.0 & 98 & 34{,}882 & 1{,}041 & 8 & 280 & 272 & 90.2 & 0 & 0 \\
 & \texttt{epoll-split} & 3/3 & 27{,}200 & 0 & 4.2 & 8.7 & 10.5 & 79 & 20{,}995 & 4{,}578 & 9 & 282 & 272 & 94.4 & 0 & 0 \\
\midrule
1{,}040 & \texttt{threads} & 3/3 & 15{,}584 & 1{,}736{,}330 & 18{,}253.6 & 18{,}580.8 & 18{,}580.8 & 192 & 89{,}495 & 2{,}866 & 1{,}032 & 2{,}054 & 1{,}024 & 310.4 & 0 & 0 \\
 & \texttt{pool-sender} & 3/3 & 13{,}590 & 1{,}776{,}200 & 18{,}790.5 & 19{,}014.6 & 19{,}014.6 & 192 & 90{,}053 & 2{,}887 & 1{,}033 & 2{,}055 & 1{,}024 & 302.0 & 0 & 0 \\
 & \texttt{peer-threads} & 3/3 & 19{,}592 & 1{,}656{,}160 & 17{,}179.9 & 17{,}827.1 & 17{,}827.1 & 190 & 107{,}437 & 3{,}886 & 1{,}032 & 2{,}054 & 1{,}024 & 155.2 & 0 & 0 \\
 & \texttt{poll} & 3/3 & 102{,}400 & 0 & 205.5 & 427.8 & 486.5 & 133 & 6{,}324 & 588 & 8 & 2{,}055 & 1{,}024 & 285.2 & 0 & 0 \\
 & \texttt{poll-duplex} & 3/3 & 102{,}400 & 0 & 71.3 & 151.0 & 209.7 & 145 & 16{,}234 & 732 & 8 & 1{,}031 & 1{,}024 & 75.5 & 0 & 0 \\
 & \texttt{poll-split} & 3/3 & 102{,}400 & 0 & 58.7 & 127.9 & 172.0 & 156 & 16{,}049 & 2{,}281 & 9 & 1{,}032 & 1{,}024 & 90.2 & 0 & 0 \\
 & \texttt{epoll} & 3/3 & 102{,}400 & 0 & 176.2 & 771.8 & 1{,}140.9 & 141 & 9{,}638 & 577 & 8 & 2{,}056 & 1{,}024 & 285.2 & 0 & 0 \\
 & \texttt{epoll-duplex} & 3/3 & 102{,}400 & 0 & 159.4 & 687.9 & 1{,}006.6 & 139 & 10{,}171 & 503 & 8 & 1{,}032 & 1{,}024 & 92.3 & 0 & 0 \\
 & \texttt{epoll-split} & 3/3 & 102{,}400 & 0 & 138.4 & 620.8 & 922.7 & 143 & 10{,}404 & 1{,}185 & 9 & 1{,}034 & 1{,}024 & 98.6 & 0 & 0 \\
\end{longtable}

\end{landscape}
